%==============================================================================
%  Section 6 --- NEPANTLA
%  BODY ONLY. The preamble lives in ../nepantla.tex
%
%  NOTE ON NUMBERING: the master shares one counter across theorem / proposition
%  / corollary. Sections 3's two propositions consume 1 and 2, so the counter is
%  reset here twice so that the results print under the names the paper uses
%  throughout -- "Theorem 1", "Corollary 1", "Corollary 2".
%==============================================================================

\section{NEPANTLA}
\label{sec:nepantla}

\noindent\emph{Non-inferior Execution via Progressive Augmentation, Native-Token
Locking and Arbitration.}

\subsection{The idea in one paragraph}
\label{sec:nepantla-idea}

Stop asking \emph{``is this model good at Spanish?''} and start asking
\emph{``is this proposed action correct?''} The first question is unanswerable
in advance for a model the user chose and we have never seen. The second is
largely \textbf{decidable}, because an operator system's decisions are
structured objects --- a tool name from a fixed vocabulary, argument keys from a
fixed schema, argument values that must trace back to something the user
actually wrote. NEPANTLA therefore freezes the user's literals before anything
can touch them, proposes an action at the cheapest rung of an escalation ladder,
submits that proposal to a language-independent verifier \textbf{before it
executes}, and climbs the ladder on rejection. The ladder's last rung is
operationally the English pipeline. Since every earlier rung is \emph{attempted,
verified and discarded} without side effects, the ladder cannot end below its
own last rung --- and the answer the user reads is rendered into Spanish
afterwards, by a channel that Proposition~1 proves cannot touch correctness at
all.

\subsection{Architecture}
\label{sec:nepantla-architecture}

\begin{figure}[htbp]
\centering
\singlespacing
\colorlet{npaccent}{blue!55!black}
\begin{tikzpicture}[
    font=\scriptsize,
    node distance=5mm,
    every node/.style={align=center},
    box/.style   ={draw=black!55, rounded corners=2pt, fill=black!5,
                   text width=0.56\textwidth, inner sep=4pt},
    stage/.style ={box, thick, fill=black!10},
    io/.style    ={draw=black!55, rounded corners=7pt, fill=black!3,
                   text width=0.50\textwidth, inner sep=4pt},
    rung/.style  ={draw=black!55, rounded corners=2pt, fill=white,
                   text width=0.38\textwidth, inner sep=3pt},
    term/.style  ={rung, thick, draw=npaccent, fill=npaccent!8},
    probe/.style ={draw=black!45, dashed, rounded corners=2pt, fill=black!3,
                   text width=0.19\textwidth, inner sep=3pt},
    flow/.style  ={-{Stealth[length=2mm]}, thick, draw=black!65},
    side/.style  ={-{Stealth[length=2mm]}, dashed, draw=black!50},
    lbl/.style   ={font=\scriptsize, inner sep=1.5pt}
  ]

  % ------------------------------------------------------------------ input
  \node[io] (U) {Utterance $u$ (Spanish)};

  % ------------------------------------------------------------- stage 1 + 2
  \node[stage, below=of U] (FREEZE)
    {\textbf{STAGE 1 $\cdot$ FREEZE} --- extract literal set $\Sigma(u)$\\
     lexical, language-independent $\cdot$ NFC-normalised, byte-exact};

  \node[stage, below=of FREEZE] (NEU)
    {\textbf{STAGE 2 $\cdot$ NEUTRALISE} --- N1 fold $\cdot$ N2 boundary-match
     $\cdot$ N3 canonical expansion\\
     \emph{identity on ASCII}};

  \node[box, below=of NEU] (PLAN)
    {Language-neutral planner \textrightarrow{} expected action class};

  % ------------------------------------------------------------- the ladder
  \node[font=\scriptsize\bfseries, below=8mm of PLAN] (LADTITLE)
    {STAGE 3 $\cdot$ PROPOSE \& VERIFY (the ladder)};

  \node[rung, below=2mm of LADTITLE] (R0)
    {\textbf{R0 $\cdot$ NATIVE}\\ Spanish verbatim + EN spine};
  \node[rung, below=4.5mm of R0] (R1)
    {\textbf{R1 $\cdot$ ANCHORED}\\ + explicit $\Sigma$ literal table};
  \node[rung, below=4.5mm of R1] (R2)
    {\textbf{R2 $\cdot$ NEPANTLA RUNG}\\ + EN gloss beside the Spanish\\
     \emph{augmentation, never substitution}};
  \node[term, below=4.5mm of R2] (R3)
    {\textbf{R3 $\cdot$ ENGLISH-EQUIVALENT}\\ = the English baseline};
  \node[rung, below=4.5mm of R3] (R4)
    {\textbf{R4 $\cdot$ HONEST STOP}\\ refuse in Spanish, name the reason};

  \node[draw=black!50, dashed, rounded corners=3pt, inner sep=5pt,
        fit=(LADTITLE)(R0)(R1)(R2)(R3)(R4)] (LADDER) {};

  % ------------------------------------------------------------- the verifier
  \node[draw=npaccent, thick, diamond, aspect=3.6, inner sep=1pt,
        fill=npaccent!10, text width=3.4cm, below=7mm of LADDER] (V)
    {\textbf{VERIFIER $V$}\\ 6 language-independent checks\\
     \emph{runs BEFORE execution}};

  % ------------------------------------------------------- execute and render
  \node[box, below=6mm of V, text width=0.44\textwidth] (EXEC)
    {\textbf{EXECUTE} --- $\tau$ \textrightarrow{} machine state};

  \node[stage, below=6mm of EXEC] (RENDER)
    {\textbf{STAGE 4 $\cdot$ RENDER} --- Spanish presentation, protected spans
     opaque\\ \emph{downstream of the decision}};

  \node[io, below=of RENDER] (OUT)
    {Spanish answer + Spanish Exec Report\\
     English asset names verbatim};

  % --------------------------------------------------------- capability probe
  \node[probe, right=6mm of R1] (PROBE)
    {capability profile\\ \emph{chooses the START rung only}};
  \node[lbl, below=1.5mm of PROBE, text width=0.19\textwidth]
    {\emph{optimisation,\\ not a gate}};

  % ------------------------------------------------------------------- edges
  \draw[flow] (U)      -- (FREEZE);
  \draw[flow] (FREEZE) -- (NEU);
  \draw[flow] (NEU)    -- (PLAN);
  \draw[flow] (PLAN)   -- (LADDER);

  \draw[flow] (R0) -- node[lbl, right] {$V$ rejects} (R1);
  \draw[flow] (R1) -- node[lbl, right] {$V$ rejects} (R2);
  \draw[flow] (R2) -- node[lbl, right] {$V$ rejects} (R3);
  \draw[flow] (R3) -- node[lbl, right] {$V$ rejects} (R4);

  \draw[side] (PROBE) -- (R1);

  \draw[flow] (LADDER) -- node[lbl, right] {proposal} (V);
  \draw[flow] (V)      -- node[lbl, right] {accept}   (EXEC);
  \draw[flow] (V.west) -- ++(-1.7,0)
              |- node[lbl, right, pos=0.4] {reject($\rho$)} (LADDER.west);

  \draw[flow] (EXEC)   -- (RENDER);
  \draw[flow] (RENDER) -- (OUT);

\end{tikzpicture}
\caption{The NEPANTLA architecture. The literal set is frozen before any
transformation touches the utterance; the deterministic core is neutralised by
operators that are the identity on ASCII; every proposed action is verified
\emph{before} it executes, so escalation up the ladder is side-effect-free; and
the Spanish presentation is rendered strictly downstream of the verified
decision. The capability profile only chooses which rung to start at.}
\label{fig:nepantla-architecture}
\end{figure}

\subsection{Stage 1 --- Native-Token Locking}
\label{sec:native-token-locking}

Before the utterance touches a model, a prompt template, a scorer or a
transformation of any kind, a \textbf{lexical extractor} walks it and produces
the frozen set $\Sigma(u)$.

The extractor is deliberately \emph{shape-based, not meaning-based} --- it
recognizes the syntactic silhouette of a literal, never its semantics --- which
is exactly why it is language-independent:

\begin{table}[htbp]
\centering
\footnotesize
\caption{The literal extractor recognizes shapes, not meanings.}
\label{tab:literal-extractor}
\begin{tabular}{@{}>{\raggedright\arraybackslash}p{0.14\textwidth}
                  >{\raggedright\arraybackslash}p{0.31\textwidth}
                  >{\raggedright\arraybackslash}p{0.37\textwidth}@{}}
\toprule
\textbf{Class} & \textbf{Recognizer (sketch)} & \textbf{Example} \\
\midrule
Windows path & drive letter, colon, backslash run &
  \code{C:\textbackslash{}Tlamatini\textbackslash{}Templates\textbackslash{}leg\_ctrl} \\
POSIX path & $\ge 2$ slash-separated segments & \code{/usr/local/bin/pio} \\
Filename & stem + known or generic extension & \code{informe\_año.pdf} \\
Glob / regex & wildcard or metacharacter density &
  \code{*.log}, \code{\textasciicircum{}ERROR:.*\$} \\
Flag & leading \code{-} / \code{--} + identifier &
  \code{--noreload}, \code{-sT} \\
Quantity & digits with optional unit/separator &
  \code{115200}, \code{1.5}, \code{8 GB} \\
Endpoint & host:port, URL, e-mail & \code{127.0.0.1:5000} \\
Machine identifier & exact match against the termbase &
  \code{Emailer}, \code{chat\_agent\_apirer}, \code{bluepill\_f103c8} \\
Quoted span & balanced quotes & \code{'hola'} \\
\bottomrule
\end{tabular}
\end{table}

Three properties matter.

\begin{enumerate}
  \item \textbf{It runs first.} Nothing precedes it. Whatever the rest of the
        pipeline does, $\Sigma$ is already the ground truth.
  \item \textbf{It over-extracts deliberately.} A false positive costs one
        unnecessary provenance entry; a false negative removes a literal from
        protection. The asymmetry dictates the tuning.
  \item \textbf{NFC only.} Composed and decomposed forms of the same character
        are unified; case and accents are \emph{never} folded, because a path
        differing by an accent is a different path (\S3.3).
\end{enumerate}

$\Sigma$ then travels with the request as immutable metadata, and it is used
three times: in the anchor table at R1, in the gloss's protected-span mask at
R2, and in the verifier's provenance check at every rung.

\subsection{Stage 3 --- The verifier \texorpdfstring{$V$}{V}}
\label{sec:nepantla-verifier}

$V$ is the heart of the guarantee. It inspects a \textbf{proposed} action before
execution and returns accept or reject-with-reason. It knows nothing about
Spanish, nothing about English, and nothing about the model.

\begin{table}[htbp]
\centering
\footnotesize
\caption{The verifier's checks. Every one of them is decidable, which is a
property of \emph{operator} systems rather than an accident of this one.}
\label{tab:verifier}
\begin{tabular}{@{}l
                  >{\raggedright\arraybackslash}p{0.33\textwidth}
                  >{\raggedright\arraybackslash}p{0.13\textwidth}
                  >{\raggedright\arraybackslash}p{0.32\textwidth}@{}}
\toprule
\textbf{\#} & \textbf{Check} & \textbf{Decidable?} & \textbf{Catches} \\
\midrule
\textbf{V1} & \textbf{Tool existence} --- $\mathrm{name} \in \mathcal{N}$, the
  bound surface, byte-exact ASCII & \textbf{Yes} &
  Hallucinated tools; a model inventing \code{enviar\_correo} \\
\addlinespace[2pt]
\textbf{V2} & \textbf{Schema conformance} --- required keys present, types
  valid, enums from the declared set & \textbf{Yes} &
  Missing arguments; a translated enum value \\
\addlinespace[2pt]
\textbf{V3} & \textbf{Argument provenance} --- every literal in \code{args}
  $\in \Sigma(u) \cup \mathcal{R}(\mathcal{H}) \cup \mathcal{D}$ & \textbf{Yes} &
  \emph{The cross-lingual failure class}: normalized paths, stripped accents,
  translated flags, re-formatted numbers, invented filenames \\
\addlinespace[2pt]
\textbf{V4} & \textbf{Precondition satisfaction} --- the target exists / does
  not exist as required; the board is attached; the port is open &
  \textbf{Yes} (via the existing fail-safe preflights) &
  Actions that cannot succeed \\
\addlinespace[2pt]
\textbf{V5} & \textbf{Gating parity} --- the Ask-Execs tier of the chosen tool
  matches the tier implied by the planner's expected action class &
  \textbf{Yes} &
  \emph{Safety regression}: a Spanish request routed to an ungated
  near-synonym \\
\addlinespace[2pt]
\textbf{V6} & \textbf{Sentinel integrity} --- every protocol token emitted is
  byte-exact & \textbf{Yes} &
  A mangled \code{INI\_SECTION\_\dots} or \code{VERDICT:} corrupting downstream
  routing \\
\addlinespace[2pt]
\textbf{V7} & \textbf{Action expectancy} --- if the planner expected an action
  class and the model returned only prose, reject & \textbf{Yes} &
  The \S5.3 failure: a destructive request answered with advice \\
\bottomrule
\end{tabular}
\end{table}

Every check is decidable. That is not an accident of this system; it is a
property of \emph{operator} systems generally, and it is what makes the
guarantee reachable here and not in a pure-prose product.

\textbf{V3 deserves its own remark}, because it is the paper's most transferable
idea. Both the Spanish and English renderings of a task contain the \emph{same
frozen literals} by construction. Therefore a check of the form \emph{``can this
emitted value be traced to something the user actually wrote, or to a prior
verified result, or to a declared default?''} is simultaneously (a) completely
language-agnostic, (b) cheap, and (c) precisely aimed at the damage
cross-lingual operation causes. A value with no provenance is either a
hallucination or a translation artifact --- and in an operator system those are
the same emergency.

\textbf{What $V$ does not check.} It cannot decide whether the
\emph{semantically most appropriate} tool was chosen among several valid ones.
That residue is exactly the part NEPANTLA inherits from the English baseline
rather than improving on, and \S6.7 is careful to claim only that.

\textbf{Rejection is cheap and side-effect-free} because $V$ runs
\emph{pre-execution}. A rejected proposal has changed nothing; escalation
replays planning, never state. This is the property that makes the ladder sound.

\subsection{Stage 3 --- The escalation ladder}
\label{sec:escalation-ladder}

Each rung differs from the previous only by \textbf{adding information}, never
by removing or replacing it.

\begin{description}

\item[R0 --- NATIVE.] English system spine (byte-stable,
prompt-cache-friendly), Spanish user content verbatim, Spanish answer directive,
low temperature. Cheapest and highest-fidelity: zero transformations, so
$p^{k} = 1$ trivially.

\emph{Reached first when the capability profile says the model is
Spanish-competent.}

\item[R1 --- ANCHORED NATIVE.] R0 plus an explicit, machine-readable anchor
table appended to the prompt:

\begin{lstlisting}[caption={The R1 anchor table, appended verbatim to the prompt.},label={lst:anchor-table}]
LITERALS --- reproduce these byte-for-byte; never translate, never normalise:
  L1 = C:\Tlamatini\Templates\leg_ctrl
  L2 = bluepill_f103c8
  L3 = 115200
\end{lstlisting}

This converts an implicit expectation into an explicit one and repairs the most
common V3 rejection at negligible cost. Empirically this is where most
Spanish-competent-but-careless models are recovered.

\item[R2 --- THE NEPANTLA RUNG (bilingual augmentation).] R1 plus an
\textbf{English gloss of the intent placed beside the Spanish original}, never
instead of it. The gloss is produced with $\Sigma$ masked out and re-injected
verbatim, so it is structurally incapable of damaging a literal.

The model now sees both renderings simultaneously. A Spanish-competent model
ignores the gloss as redundant; a Spanish-weak model reads the gloss and
recovers the intent. This is the rung the system's name refers to: the middle
place, where both languages are present and neither is authoritative over the
literals.

\emph{Note the placement.} The gloss is a translation, and it sits
\textbf{upstream} of a decision --- which Principle~T forbids. The contradiction
is only apparent: Principle~T forbids upstream translation \emph{on the causal
path to the literals}. Here the literals bypass the gloss entirely (they are
masked, then re-injected, then verified by V3), so the only thing the
translation can influence is \emph{which tool} is chosen --- and that influence
is additive evidence, subject to the same verifier.

\item[R3 --- ENGLISH-EQUIVALENT EXECUTION.] The execution decision is driven by
the English rendering under exactly the baseline's system prompt, tool schemas
and parameters. This rung is, operationally, $B$.

Crucially, \textbf{the user still receives Spanish}, because the answer is
produced by Stage~4, which by Proposition~1 is causally downstream of $\tau$.
Reaching R3 is invisible to the operator except as a small latency cost and a
status note.

\item[R4 --- HONEST STOP.] If even R3's proposal fails verification, NEPANTLA
does not guess. It refuses, in Spanish, naming the verifier's rejection reason
and the literal or tool involved, and it suggests the concrete remedy (rephrase,
supply the missing path, or switch model). No action is executed.

This rung exists because the only thing worse than an English answer to a
Spanish operator is a \textbf{confident wrong action}. Refusal is a correct
outcome; fabrication is not.

\end{description}

\subsection{Stage 4 --- The presentation renderer}
\label{sec:presentation-renderer}

Independent channel, no authority over $\tau$, three ordered strategies:

\begin{enumerate}
  \item \textbf{Native.} The model already answered in Spanish (the usual case,
        including at R3 when the model is competent). Nothing to do.
  \item \textbf{Delegated verbalization.} The decisions and results are fixed
        and verified; a \emph{separate}, known-Spanish-capable renderer (a small
        local model, or the static catalog for structured content) turns them
        into Spanish prose. This is downstream translation, licensed by
        Principle~T, and it cannot alter a single byte of what was executed.
  \item \textbf{Structured fallback.} If no verbalizer is available, the answer
        is composed from the \textbf{Spanish message catalog} plus the verified
        structured results --- a fully Spanish, slightly terser answer assembled
        by template rather than generated. The Exec Report already provides most
        of what the operator needs, and its chrome is catalog-driven.
\end{enumerate}

In all three, \textbf{protected spans are opaque}: agent names, tool names,
paths, flags, code blocks, sentinels and the brand pass through byte-identical.
\S2.4's rendering table is the specification.

A \textbf{read-only guard} runs over the raw model output, before any stripping,
and reports two things to the log (\code{--- [I18N-GUARD]}) without ever
mutating the answer: sentinel integrity, and answer-language line-pass-rate over
\emph{masked prose only}. It is read-only by policy: an answer-rewriting pass
driven by an uncalibrated language detector is a corruption engine, and because
the response pipeline persists the answer, a false repair would be replayed on
every chat reload forever. Repair is unlocked only by a measured false-positive
rate.

\subsection{The algorithm}
\label{sec:nepantla-algorithm}

\begin{lstlisting}[language=Python,caption={NEPANTLA --- the complete algorithm. The verification loop sits between proposal and execution.},label={lst:nepantla}]
def NEPANTLA(u, session, model, profile):
    # ---- STAGE 1: freeze. Nothing has touched u yet. -----------------------
    Sigma  = extract_literals(u)                 # lexical, language-independent
    lang   = resolve_language(u, session)        # conversation-sticky, hysteretic
    H      = session.verified_results            # provenance from prior turns

    # ---- STAGE 2: neutralise the deterministic core ------------------------
    key    = N3(N1(u), lang)                     # fold + canonical expansion
    plan   = planner(key, phrase_match=N2)       # expected action class, tools
    # identity on ASCII: an English u yields today's exact plan

    # ---- STAGE 3: propose -> verify -> escalate ----------------------------
    start  = profile.start_rung(model, lang)     # OPTIMISATION ONLY (Cor. 1)
    trace, notes = [], []

    for rung in LADDER[start:]:                  # R0, R1, R2, R3, R4
        if rung is R4:
            return honest_stop(lang, notes, Sigma)      # refuse, in Spanish

        prompt   = build_prompt(rung, u, Sigma, plan, lang, model)
        proposal = model.propose(prompt)         # actions are NOT executed yet

        verdicts = [V(Sigma, H, a) for a in proposal.actions]
        if all(v.accepted for v in verdicts):
            trace = proposal.actions
            break

        notes.append(rejection_summary(verdicts))
        # no side effects occurred: escalation replays planning, never state

    # ---- EXECUTE the verified trace ---------------------------------------
    results = []
    for a in trace:
        r = execute(a)                           # existing permission gate applies
        if failed(r):                            # post-hoc failure (preconditions
            notes.append(r)                      # can change between check and run)
            corrective_loop(a, r)                # the system's existing machinery
        else:
            H.record(a, r)                       # literals enter provenance
            results.append(r)

    # ---- STAGE 4: render, downstream of the decision -----------------------
    answer = render(results, lang, protected=Sigma | TERMBASE, notes=notes)
    guard_report = inspect_readonly(model.raw_output, lang)   # never mutates
    log(guard_report)
    return answer
\end{lstlisting}

Two lines carry most of the weight. \code{proposal = model.propose(prompt)}
\textbf{does not execute}, and the verification loop sits between proposal and
execution. Everything else is bookkeeping around that separation.

\subsection{Theorem 1 --- Ladder Dominance}
\label{sec:ladder-dominance}

\setcounter{theorem}{0}
\begin{theorem}[Ladder Dominance]
\label{thm:dominance}
Let the ladder be $R_0, \dots, R_{m}$ with terminal executable rung $R_{m-1}$
and honest stop $R_m$. Let $\mathrm{FA}_i$ be the probability that $V$
\textbf{falsely accepts} an incorrect proposal at rung $i$. Then
\begin{align*}
  \theta_{A} \;&\ge\; \theta_{R_{m-1}} \;-\; \sum_{i<m-1}\mathrm{FA}_i ,\\
  \intertext{and if additionally $\theta_{R_{m-1}} \ge \theta_{EN}$, then}
  \theta_{A} \;&\ge\; \theta_{EN} \;-\; \mathrm{FA}_{\text{total}} .
\end{align*}
\end{theorem}

\begin{proof}
Consider any task on which the terminal rung would succeed. NEPANTLA fails on
that task only if it terminated \emph{before} reaching $R_{m-1}$ with an
incorrect trace. Termination before $R_{m-1}$ occurs only on acceptance by $V$.
Therefore the event ``$A$ fails while $R_{m-1}$ would have succeeded'' is
contained in the union over earlier rungs of ``$V$ falsely accepted an incorrect
proposal'', whose probability is at most $\sum_{i<m-1}\mathrm{FA}_i$ by a union
bound. Rejection at an earlier rung is side-effect-free by the pre-execution
property of \S6.4, so it cannot itself cause failure --- it can only cost
latency. Hence the first inequality; the second follows by substitution.
\end{proof}

\textbf{Why $\mathrm{FA}$ is zero on the failure classes that matter.} Checks
V1, V2, V3, V6 and V7 are \emph{decidable predicates over the proposal}, so
their false-accept probability is exactly $0$: a nonexistent tool, a schema
violation, a literal with no provenance, a damaged sentinel and a
missing-but-expected action are each detected with certainty. By
Proposition~2, literal corruption in particular has $\mathrm{FA}=0$. The
residual $\mathrm{FA}$ is confined to V4 (preconditions can change between check
and execution --- a genuine time-of-check/time-of-use window) and to the
\emph{undecidable} residue $V$ deliberately does not attempt: choosing the
semantically best tool among several schema-valid candidates.

That residue is the honest boundary of the claim, and it is worth stating in
plain language:

\begin{keyinsight}
\textbf{NEPANTLA is non-inferior to the English pipeline on every failure class
that operating in Spanish actually introduces. On the residual class --- which
tool is the wisest choice --- it inherits the English pipeline's behaviour
rather than improving on it.}
\end{keyinsight}

\textbf{Discharging the condition $\theta_{R_{m-1}} \ge \theta_{EN}$.} Three
arguments support it and one caveat bounds it.

\begin{enumerate}[label=(\roman*), leftmargin=*]
  \item The literals reaching $R_3$ are byte-identical to the user's own, by
        Proposition~2 --- so on the argument axis $R_3$ equals the baseline
        exactly.
  \item $R_3$'s system spine, tool schemas and generation parameters are the
        baseline's, unmodified.
  \item $R_3$ presents \emph{both} the English rendering and the Spanish
        original, so the model's evidence is a superset of what a gloss alone
        would give.
\end{enumerate}

\emph{Caveat:} additional context is not provably monotone for a language model
--- more text can distract as well as inform. The condition is therefore an
\textbf{empirical} one, and \S8 exists to test precisely it. A paper that
claimed otherwise would be overselling; the architecture guarantees the
\emph{ladder} property unconditionally, and measurement establishes the
\emph{terminal} property.

\textbf{False rejections.} If $V$ wrongly rejects a correct proposal, the cost
is an unnecessary escalation --- latency --- and in the worst case an honest
refusal at R4. The failure direction is toward \emph{not acting}, which for a
system that deletes files and flashes firmware is the correct direction to fail
in.

\subsection{Corollary 1 --- Probe Safety}
\label{sec:probe-safety}

\setcounter{theorem}{0}
\begin{corollary}[Probe Safety]
\label{cor:probe-safety}
The model-capability profile affects only the expected number of rungs
evaluated. It cannot affect the correctness of the terminal outcome.
\end{corollary}

\begin{proof}
The profile appears in the algorithm exactly once, as \code{start} --- an index
into the ladder. Theorem~\ref{thm:dominance}'s bound is over the suffix
beginning at any start index, since every rung after \code{start} remains
reachable and the terminal rung is always the same.
\end{proof}

This has a strong practical consequence, and it is the answer to the obvious
objection \emph{``what if you mis-measure the model?''}:

\begin{keyinsight}
\textbf{A miscalibrated capability probe costs latency, never correctness.}
\end{keyinsight}

The profile is therefore permitted to be a cheap, imperfect heuristic. It is an
accelerator, not a gate --- which in turn means it does not need the elaborate
calibration machinery a \emph{gating} probe would demand, and it can never
silently deny a user their language because a cache file was unreadable.

\S7.3 shows that this permission is not a convenience but a
\textbf{necessity}: no provider exposes the information a gating probe would
need, so a design in which correctness depended on measuring the model
correctly would be undeliverable. Corollary~\ref{cor:probe-safety} is what
makes the measurement problem of \S7.3 tractable at all.

\textbf{Amendment --- the profile appears twice.} The proof above is exact for
the use it describes, and it is the only use in the pseudocode of \S6.7, where
Stage~2 reads \code{key = N3(N1(u), lang)} --- gated on the language alone. The
\textbf{implemented} Stage~0 (\code{agent/i18n/model\_caps.py}) gates N3's
expansion on the model's measured tier as well, so the profile also influences
the deterministic scorer \emph{upstream} of the ladder. That second appearance
is not covered by Corollary~\ref{cor:probe-safety}, and it would be dishonest
to let the corollary's blanket phrasing absorb it. It is covered instead by a
second, deliberately weaker result.

% The prime belongs to the result's NAME ("Corollary 1'"), not to the shared
% counter. \thetheorem (and hyperref's \theHtheorem, so the PDF anchor stays a
% plain string) are redefined inside a group so that the printed header and
% \ref{cor:assist-safety} agree; the counter is then restored to 1 so that
% "Corollary 2" below still prints as 2.
\begingroup
\renewcommand{\thetheorem}{1\ensuremath{'}}
\def\theHtheorem{1prime}
\begin{corollary}[Assist Safety]
\label{cor:assist-safety}
A mis-tiered model can perturb the planner's \emph{ranking} and therefore the
expected number of rungs and the bind order under context overflow. It cannot
alter the action vocabulary $\mathcal{N}$, cannot introduce a capability that
does not exist, and cannot change the terminal rung.
\end{corollary}
\endgroup
\setcounter{theorem}{1}

\begin{proof}
N3 appends only hint tokens \textbf{that already exist in the registry},
enforced by the closure test of \S5.9. Expansion can therefore raise the score
of an existing capability but can never invent one, and it leaves $\mathcal{N}$
pointwise unchanged; N1 is the identity on ASCII and N2 only removes matches.
The planner's output is a ranking and an expected action class --- not an
action. Every candidate action still passes $V$ before execution, and $R_3$
remains reachable with the baseline's own planning, so the terminal rung is
unaffected and Theorem~\ref{thm:dominance}'s suffix argument applies unchanged.
\end{proof}

Unlike Corollary~\ref{cor:probe-safety}, this is \textbf{not} a zero-impact
claim, and the exposure should be named. When the full tool surface exceeds the
context budget, bind order is decided by the same scorer (\S5.3), and a tool
that is never bound can never be proposed --- which is a failure $V$ cannot
repair, because $V$ only inspects proposals it is given. Two facts bound that
exposure. Multi-Turn binds the \textbf{full} enabled surface whenever it fits,
so overflow is the only window in which ranking becomes selection. And the
fail-open direction of \S7.3 points toward \emph{more} canonical keys rather
than fewer, which is the direction that helps a genuinely relevant capability
cross the threshold rather than the direction that hides one. The residual risk
is therefore contained inside Theorem~\ref{thm:dominance}'s already-declared
undecidable residue --- \emph{which schema-valid tool is wisest} --- and does
not open a new failure class.

\subsection{Corollary 2 --- Naming Necessity}
\label{sec:naming-necessity}

\begin{corollary}[Naming Necessity]
\label{cor:naming}
If any element of the action vocabulary $\mathcal{N}$ --- a tool name, an agent
display name, an argument key, an enumerated value, a protocol sentinel --- is
localized, then Proposition~1 fails, $V$'s checks V1, V2 and V6 become
locale-relative, and Theorem~\ref{thm:dominance} no longer applies.
\end{corollary}

\begin{proof}
V1 tests $\mathrm{name} \in \mathcal{N}$. If $\mathcal{N}$ is indexed by locale,
the test is $\mathrm{name} \in \mathcal{N}_\ell$, and the English baseline's
action space $\mathcal{N}_{en}$ and the Spanish pipeline's $\mathcal{N}_{es}$
are different sets. The comparison $\theta_{ES} \ge \theta_{EN}$ then ranges
over incomparable outcome spaces, and $\mathrm{exec}$ acquires a dependence on
$\ell(u)$, contradicting Proposition~1.
\end{proof}

In engineering terms, and this is the sentence to remember:

\begin{keyinsight}
\textbf{\code{Emailer} must stay \code{Emailer} in the Spanish build not because
translating it would be untidy, but because translating it deletes the proof.}
\end{keyinsight}

The same holds for \code{Asker}, \code{Apirer}, \code{STM32er},
\code{Kyber-KeyGen}, \code{File-Creator}, every tool name, every flag, every
configuration key, every sentinel, and every identifier in the source. The
\S2.3 table is not a style guide; it is the hypothesis of Proposition~1 written
out longhand.

\textbf{The escape route, and why it is closed.} Corollary~\ref{cor:naming} as
proved establishes a \emph{conditional} necessity: localizing the vocabulary
breaks \textbf{this} proof. That leaves an objection standing, and it is the
strongest one available against \S2.3. An objector may accept the corollary
entirely and reply:

\begin{quote}
``Then do not use one vocabulary. Index it. Keep a per-locale registry in which
\code{Correero} is the Spanish name of the agent whose English name is
\code{Emailer}, resolve the index at bind time, and you obtain a system with a
localized surface and an equally sound proof --- a proof over
$\mathcal{N}_\ell$ rather than over $\mathcal{N}$.''
\end{quote}

That construction is not refuted by Proposition~1. It is a different
architecture, and on paper it is coherent. It requires exactly one thing: an
\textbf{oracle that says which index applies} --- which vocabulary a given bound
model expects, or equivalently which language that model is operating in with
sufficient reliability to route on. Without such an oracle the resolution step
is a guess, and a guess about \emph{which name set the model will emit from} is
upstream of V1, which is precisely the position Principle~T forbids.

\S7.3 establishes, by verification against five providers' primary schemas,
that \textbf{no such oracle is published by anyone}, and that the single
structured language declaration that exists anywhere in the stack is absent
from half the multilingual models that carry it, wrong where present, and
inverted by the transport that serves it. The escape route therefore has no
input.

The upgrade this delivers should be stated precisely, because overstating it
would be the same error the paper criticizes elsewhere. We have \textbf{not}
proved that a locale-indexed action vocabulary is impossible in principle; a
provider could publish the required metadata tomorrow, and the argument would
have to be reopened. What is established is narrower and, for an engineering
decision, sufficient:

\begin{keyinsight}
\textbf{Corollary~\ref{cor:naming} is necessary for this proof. The only
construction that would evade it is unimplementable on the information any
provider actually publishes.}
\end{keyinsight}

The evidence of \S7.3 therefore \textbf{strengthens} Corollary~\ref{cor:naming},
and it does so along an axis the original derivation could not reach: the
original argues from inside the framework, the measurement argues from outside
it and closes the alternative. One finding discharges two obligations --- the
same absence of a language oracle that \emph{forces} \S7.3 to measure rather
than look up is what \emph{forbids} the locale-indexed vocabulary that would
otherwise let \S2.3 be softened. Naming invariance and capability measurement
are not two independent design choices. They are the two consequences of a
single empirical fact.

\subsection{Complexity and cost}
\label{sec:complexity-and-cost}

Let $q_i$ be the probability that the proposal at rung $i$ is rejected. Expected
model calls per request:
\[
  \mathbb{E}[\text{calls}] \;=\; 1 + q_{0} + q_{0}q_{1} + q_{0}q_{1}q_{2}
\]
bounded above by 4 and, for a Spanish-competent model where $q_0$ is small,
approximately $1 + q_0$. With the profile starting a competent model at R0 and
an unknown model at R1, the measured cost target is \textbf{under 1.15 model
calls per request} --- that is, the Spanish path costs essentially what English
costs, and pays extra only exactly when it was about to be wrong.

The non-model costs are all $O(n)$ in the utterance length and negligible
against a model call:

\begin{table}[htbp]
\centering
\small
\caption{Non-model cost of the language layer, per request.}
\label{tab:nepantla-cost}
\begin{tabular}{@{}>{\raggedright\arraybackslash}p{0.24\textwidth}
                  >{\raggedright\arraybackslash}p{0.42\textwidth}
                  >{\raggedright\arraybackslash}p{0.18\textwidth}@{}}
\toprule
\textbf{Stage} & \textbf{Cost} & \textbf{Budget} \\
\midrule
Literal extraction & one linear pass, compiled patterns & $<1$\,ms \\
N1/N2/N3 & one fold, cached boundary patterns, dict lookups & $<1$\,ms \\
Verifier per action & set membership + schema walk + preflight &
  $<2$\,ms (preflight dominates) \\
Renderer protected-span masking & linear, short-circuits entirely when the
  answer is already Spanish & $<1$\,ms \\
\bottomrule
\end{tabular}
\end{table}

Two design rules keep this true. The English path short-circuits the entire
language layer at the first comparison, so an English user pays a single branch.
And the capability profile is computed \textbf{out of band} --- never on the
request path, one global worker, refusing to run while a generation is in
flight, because on a single local GPU ``off the call stack'' is not the same as
``off the resource''.

\subsection{Termination, cancellation and side-effect safety}
\label{sec:termination-cancellation}

\textbf{Termination} is structural: the ladder is a finite list and each
iteration advances the index unconditionally. There is no retry-in-place.

\textbf{Cancellation} is checked at every rung boundary and again after every
model call returns, not merely at entry. The renderer and the guard have a hard
wall-clock cap; they are polish on an answer the user has already earned, and
they must never be the reason a cancel does not take effect.

\textbf{Side-effect safety} rests on the pre-execution property: escalation
happens only over \emph{proposals}. Once execution begins, the ladder is
finished, and post-execution failures are handled by the system's existing
corrective-feedback machinery rather than by re-entering the ladder --- which
prevents the pathological case of an action being executed twice under two
different rungs.

\textbf{Idempotence at the boundary} deserves an explicit note. V4's
preconditions are checked before execution, but the world can change in between.
NEPANTLA does not attempt to close that window (doing so would require
transactional semantics over the filesystem and attached hardware); it bounds it
by checking as late as possible and by preserving the existing permission gate,
which puts a human in the loop for exactly the tier of actions where the window
matters.

\subsection{What is new here}
\label{sec:what-is-new}

The individual materials are known. The arrangement is what we claim.

\begin{itemize}
  \item \textbf{Verification-first cross-lingual operation.} The cross-lingual
        literature is dominated by \emph{making the model better} at the target
        language --- translation, scaffolds, fine-tuning, prompt engineering.
        NEPANTLA instead treats target-language competence as an unreliable
        input and moves the correctness burden onto a verifier that does not
        read either language. This reframing is what removes the dependency on a
        user-chosen model.
  \item \textbf{Argument provenance as a language-independent correctness
        certificate.} Checking that every emitted literal traces back to the
        user's own words is cheap, decidable, and exactly aligned with the
        damage cross-lingual operation causes.
  \item \textbf{The ladder-dominance argument.} A structured escalation whose
        terminal rung is the incumbent baseline converts a hard question
        (``is Spanish as good?'') into a much easier one (``is the verifier's
        false-accept rate small on the decidable classes?'').
  \item \textbf{Probe demotion.} Making capability measurement an accelerator
        rather than a gate removes an entire class of calibration risk, and is
        only possible \emph{because} of the ladder.
  \item \textbf{Naming invariance derived, not asserted.} The rule that assets
        stay English falls out of Proposition~1 as a necessary condition rather
        than being imposed as a convention --- which makes it enforceable by
        test rather than by discipline.
\end{itemize}
